% GENERATED FILE. Edit canonical lessons, then run npm run generate:books.
% canonical-source-hash: fnv1a64:536644b19c85dfc1
% canonical-lessons: SA-C194-paritosikam, SA-C194-nimantranam, SA-C194-sabha, SA-C194-gosthi, SA-C194-melakah

\chapter{Prize, Invitation, Assembly, Discussion group, Fair}
\label{ch:sa-prize-invitation-assembly-discussion-group-fair}
\hlchaptermodality{194}

\begin{chapteropening}
\textbf{By the end of this chapter:} \emph{I can say a prize, an invitation, an assembly, a discussion group and a fair.}

Complete the last lesson of chapter 194: I can say a prize, an invitation, an assembly, a discussion group and a fair.
\end{chapteropening}

\section[\texorpdfstring{\sk{पारितोषिकम्} (pāritoṣikam)}{पारितोषिकम् (pāritoṣikam)}]{\sk{पारितोषिकम्} (pāritoṣikam) — a prize}

\label{lesson:SA-C194-paritosikam}



\begin{quote}
Before the new one: say the Sanskrit for an item, then the Sanskrit for a tool.
\end{quote}

\subsection*{You'll want to know: \sk{पारितोषिकम्}}
\textbf{\sk{पारितोषिकम्}} — \emph{pāritoṣikam} — \textquotedblleft{}a prize\textquotedblright{}.

\begin{grammarlens}[title={What you've built}]
One more word for this chapter.
\end{grammarlens}

\subsection*{Your turn}
\begin{itemize}
\raggedright
  \item \emph{Say it:} \emph{pāritoṣikam}
  \item \emph{Say it:} \emph{pāritoṣikam}, once more
  \item \emph{Say it:} say \emph{pāritoṣikam}
  \item \emph{Recall:} say the Sanskrit for an item, then the Sanskrit for a tool, then say \emph{pāritoṣikam} again
\end{itemize}

\begin{tcolorbox}[breakable,colback=teal!4,colframe=teal!35!black,title={Before you move on}]
What does \sk{पारितोषिकम्} mean? (\textquotedblleft{}A prize\textquotedblright{}.) Say it once more.
\end{tcolorbox}
\section[\texorpdfstring{\sk{निमन्त्रणम्} (nimantraṇam)}{निमन्त्रणम् (nimantraṇam)}]{\sk{निमन्त्रणम्} (nimantraṇam) — an invitation}

\label{lesson:SA-C194-nimantranam}



\begin{quote}
Before the new one: say the Sanskrit for a tool, then the Sanskrit for a prize.
\end{quote}

\subsection*{You'll want to know: \sk{निमन्त्रणम्}}
\textbf{\sk{निमन्त्रणम्}} — \emph{nimantraṇam} — \textquotedblleft{}an invitation\textquotedblright{}.

\begin{grammarlens}[title={What you've built}]
One more word for this chapter.
\end{grammarlens}

\subsection*{Your turn}
\begin{itemize}
\raggedright
  \item \emph{Say it:} \emph{nimantraṇam}
  \item \emph{Say it:} \emph{nimantraṇam}, once more
  \item \emph{Say it:} say \emph{nimantraṇam}
  \item \emph{Recall:} say the Sanskrit for a tool, then the Sanskrit for a prize, then say \emph{nimantraṇam} again
\end{itemize}

\begin{tcolorbox}[breakable,colback=teal!4,colframe=teal!35!black,title={Before you move on}]
What does \sk{निमन्त्रणम्} mean? (\textquotedblleft{}An invitation\textquotedblright{}.) Say it once more.
\end{tcolorbox}
\section[\texorpdfstring{\sk{सभा} (sabhā)}{सभा (sabhā)}]{\sk{सभा} (sabhā) — an assembly, a meeting}

\label{lesson:SA-C194-sabha}



\begin{quote}
Before the new one: say the Sanskrit for a prize, then the Sanskrit for an invitation.
\end{quote}

\subsection*{You'll want to know: \sk{सभा}}
\textbf{\sk{सभा}} — \emph{sabhā} — \textquotedblleft{}an assembly, a meeting\textquotedblright{}.

\begin{grammarlens}[title={What you've built}]
One more word for this chapter.
\end{grammarlens}

\subsection*{Your turn}
\begin{itemize}
\raggedright
  \item \emph{Say it:} \emph{sabhā}
  \item \emph{Say it:} \emph{sabhā}, once more
  \item \emph{Say it:} say \emph{sabhā}
  \item \emph{Recall:} say the Sanskrit for a prize, then the Sanskrit for an invitation, then say \emph{sabhā} again
\end{itemize}

\begin{tcolorbox}[breakable,colback=teal!4,colframe=teal!35!black,title={Before you move on}]
What does \sk{सभा} mean? (\textquotedblleft{}An assembly, a meeting\textquotedblright{}.) Say it once more.
\end{tcolorbox}
\section[\texorpdfstring{\sk{गोष्ठी} (goṣṭhī)}{गोष्ठी (goṣṭhī)}]{\sk{गोष्ठी} (goṣṭhī) — a discussion group}

\label{lesson:SA-C194-gosthi}



\begin{quote}
Before the new one: say the Sanskrit for an invitation, then the Sanskrit for an assembly.
\end{quote}

\subsection*{You'll want to know: \sk{गोष्ठी}}
\textbf{\sk{गोष्ठी}} — \emph{goṣṭhī} — \textquotedblleft{}a discussion group\textquotedblright{}.

\begin{grammarlens}[title={What you've built}]
One more word for this chapter.
\end{grammarlens}

\subsection*{Your turn}
\begin{itemize}
\raggedright
  \item \emph{Say it:} \emph{goṣṭhī}
  \item \emph{Say it:} \emph{goṣṭhī}, once more
  \item \emph{Say it:} say \emph{goṣṭhī}
  \item \emph{Recall:} say the Sanskrit for an invitation, then the Sanskrit for an assembly, then say \emph{goṣṭhī} again
\end{itemize}

\begin{tcolorbox}[breakable,colback=teal!4,colframe=teal!35!black,title={Before you move on}]
What does \sk{गोष्ठी} mean? (\textquotedblleft{}A discussion group\textquotedblright{}.) Say it once more.
\end{tcolorbox}
\section[\texorpdfstring{\sk{मेलकः} (melakaḥ)}{मेलकः (melakaḥ)}]{\sk{मेलकः} (melakaḥ) — a fair, a gathering}

\label{lesson:SA-C194-melakah}



\begin{quote}
Before the new one: say the Sanskrit for an assembly, then the Sanskrit for a discussion group.
\end{quote}

\subsection*{You'll want to know: \sk{मेलकः}}
\textbf{\sk{मेलकः}} — \emph{melakaḥ} — \textquotedblleft{}a fair, a gathering\textquotedblright{}.

\begin{grammarlens}[title={What you've built}]
One more word for this chapter.
\end{grammarlens}

\subsection*{Your turn}
\begin{itemize}
\raggedright
  \item \emph{Say it:} \emph{melakaḥ}
  \item \emph{Say it:} \emph{melakaḥ}, once more
  \item \emph{Say it:} say \emph{melakaḥ}
  \item \emph{Recall:} say the Sanskrit for an assembly, then the Sanskrit for a discussion group, then say \emph{melakaḥ} again
\end{itemize}

\begin{tcolorbox}[breakable,colback=teal!4,colframe=teal!35!black,title={Before you move on}]
What does \sk{मेलकः} mean? (\textquotedblleft{}A fair, a gathering\textquotedblright{}.) Say it once more.
\end{tcolorbox}
